\documentclass[a4paper,fontset=none]{ctexart}

\usepackage{anyfontsize}
\usepackage{amsmath,amssymb,mathrsfs,wasysym}
\usepackage{autobreak}
\allowdisplaybreaks
\usepackage[T1]{fontenc}
\usepackage{textcomp}
\usepackage[libertine,vvarbb]{newtxmath}
\usepackage[scr=rsfso,frak=euler,bb=ams]{mathalfa}
\usepackage{bm}
\usepackage{indentfirst}
\setlength{\parindent}{0em}
\usepackage{parskip}
\usepackage{geometry}
\geometry{top=3cm,bottom=3cm,left=2cm,right=2cm}
\usepackage{fontspec}
\setmainfont{Libertinus Serif}
\setsansfont{Libertinus Sans}
\setmonofont{JetBrains Mono}
\setCJKmainfont{SimSun}[BoldFont=Noto Sans CJK SC Medium]

\begin{document}

    \title{\textbf{原子物理学作业}}
    \author{1900011604\ \ 张植竣\ \ 北京大学}
    \date{2021年3月24日}

    \maketitle

    \begin{itemize}

        \subsection*{思考题}

        \item[1.27]


        \item[1.32(2)]


        \item[1.33]


        \item[2.2]
        因为氨分子的翻转是分子转动能级的跃迁，苯分子共振是电子能级间的跃迁。分子的质量大约是电子的$10^4$倍，则分子翻转相应地会比电子跃迁更难，需要的能量大约也相差$10^4$倍。

        \subsection*{习题}

        \item[1.39]


        \item[2.1]
        $\bm{H}$的本征多项式为：
        \begin{align*}
        |\lambda\bm{I} - \bm{H}|
        &=
        \begin{vmatrix}
            \ \lambda-E_0 & A \\
            A & \lambda-E_0 \ \\
        \end{vmatrix} \\
        &= (\lambda - E_0)^2 - A^2 \\
        &= 0
        \end{align*}
        解得两个本征值：
        \[
        \lambda_1 = E_0 + A,\qquad \lambda_2 = E_0 - A
        \]
        对于本征值$\lambda_1 = E_0 + A$，解方程$(\lambda_1\bm{I} - \bm{H})\bm{X} = 0$：
        \[
        (\lambda_1\bm{I} - \bm{H})\bm{X}
        =
        \begin{bmatrix}
            A & A \\
            A & A \\
        \end{bmatrix}
        \begin{bmatrix}
            x_1 \\
            x_2 \\
        \end{bmatrix}
        =
        \begin{bmatrix}
            0 \\
            0 \\
        \end{bmatrix}
        \]
        解得$x_1 = -x_2$，则该方程的一个基础解系为：
        \[
        \bm{X}
        =
        \begin{bmatrix}
            1 \\
            -1 \\
        \end{bmatrix}
        \]
        对于本征值$\lambda_2 = E_0 - A$，解方程$(\lambda_1\bm{I} - \bm{H})\bm{Y} = 0$：
        \[
        (\lambda_1\bm{I} - \bm{H})\bm{Y}
        =
        \begin{bmatrix}
            -A & A \\
            A & -A \\
        \end{bmatrix}
        \begin{bmatrix}
            y_1 \\
            y_2 \\
        \end{bmatrix}
        =
        \begin{bmatrix}
            0 \\
            0 \\
        \end{bmatrix}
        \]
        解得$y_1 = y_2$，则该方程的一个基础解系为：
        \[
        \bm{Y}
        =
        \begin{bmatrix}
            1 \\
            1 \\
        \end{bmatrix}
        \]
        将$\{\bm{X},\bm{Y}\}$正交单位化得$\{\bm{\eta_1},\bm{\eta_2}\}$：
        \[
        \bm{\eta_1}
        =
        \begin{bmatrix}
            \frac{\sqrt{2}}{2} \\
            -\frac{\sqrt{2}}{2} \\
        \end{bmatrix}
        ,\qquad
        \bm{\eta_2}
        =
        \begin{bmatrix}
            \frac{\sqrt{2}}{2} \\
            \frac{\sqrt{2}}{2} \\
        \end{bmatrix}
        \]
        这就是分别属于本征值$\lambda_1 = E_0 + A$和$\lambda_2 = E_0 - A$的两个本征矢。

        对角化变换如下：

        令：
        \[
        \bm{T} = [\bm{\eta_1}, \bm{\eta_2}]
        =
        \begin{bmatrix}
            \frac{\sqrt{2}}{2} & \frac{\sqrt{2}}{2} \\
            -\frac{\sqrt{2}}{2} & \frac{\sqrt{2}}{2} \\
        \end{bmatrix}
        \]
        则：
        \[
        \bm{T}^{-1}\bm{H}\bm{T}
        =
        \begin{bmatrix}
            E_0 + A & 0 \\
            0 & E_0 - A \\
        \end{bmatrix}
        \]

        \item[3.3]


        \subsection*{补充题}

        \item[1.]
        电子中微子行进$L$距离后处于$\mu$子中微子的状态的概率为：
        \[
        P(t) = \sin^2{2\theta}\sin^2{\frac{\Delta m^2 cL}{4p\hbar}} = \sin^2{2\theta}\sin^2{\frac{\Delta E^2 L}{4Ec\hbar}}
        \]
        概率最大要求：
        \[
        \frac{\Delta E^2 L}{4Ec\hbar} = \frac{\pi}{2} + n\pi,\qquad n \in \mathbb{N}
        \]
        代入$E = 10\,\mathrm{MeV}$，$\Delta E = 1\,\mathrm{MeV}$得：
        \[
        L = (12.4 + 24.8n)\,\mathrm{m},\qquad n \in \mathbb{N}
        \]

        \item[2.]

        \item[3.]

    \end{itemize}

\end{document}